\documentclass[12pt,a4paper]{report}

\usepackage{fontspec}
\usepackage{polyglossia}
\setmainlanguage{romanian}
\setmainfont{Latin Modern Roman}

\usepackage[a4paper,margin=2.5cm]{geometry}
\usepackage{setspace}
\usepackage{amsmath}
\usepackage{amssymb}
\usepackage{enumitem}
\usepackage{longtable}
\usepackage{array}
\usepackage{booktabs}
\usepackage{xcolor}
\usepackage{listings}
\usepackage[hidelinks]{hyperref}

\onehalfspacing
\setlength{\parindent}{1.2cm}
\setlength{\parskip}{0.25em}

\definecolor{codebg}{RGB}{248,248,248}
\definecolor{codeframe}{RGB}{210,210,210}
\definecolor{codekeyword}{RGB}{0,78,130}
\definecolor{codestring}{RGB}{120,70,0}
\definecolor{codecomment}{RGB}{80,120,80}

\lstdefinestyle{pythonstyle}{
    language=Python,
    basicstyle=\ttfamily\small,
    keywordstyle=\color{codekeyword}\bfseries,
    stringstyle=\color{codestring},
    commentstyle=\color{codecomment}\itshape,
    backgroundcolor=\color{codebg},
    frame=single,
    rulecolor=\color{codeframe},
    numbers=left,
    numberstyle=\tiny,
    numbersep=8pt,
    breaklines=true,
    showstringspaces=false,
    tabsize=4,
    keepspaces=true
}

\newcommand{\codechunk}[2]{
    \lstinputlisting[
        style=pythonstyle,
        firstline=#1,
        lastline=#2,
        firstnumber=#1
    ]{../src/dsp/wiener_filter.py}
}

\title{Explicatie detaliata a fisierului \texttt{src/dsp/wiener\_filter.py}}
\author{Documentatie tehnica pentru proiectul de procesare audio}
\date{\today}

\begin{document}

\maketitle
\tableofcontents

\chapter{Rolul fisierului in proiect}

Fisierul \texttt{src/dsp/wiener\_filter.py} implementeaza o metoda clasica de reducere a zgomotului: filtrarea Wiener in domeniul timp-frecventa. In arhitectura proiectului, acest fisier apartine stratului DSP, adica stratului de metode deterministe, bazate pe semnal, care pot fi folosite ca alternative sau ca repere fata de metoda neuronala bazata pe masca spectrala.

Din punct de vedere functional, fisierul face trei lucruri principale:

\begin{itemize}
    \item estimeaza densitatea spectrala de putere a zgomotului, pornind de la cadrele cu energie mica;
    \item aplica un castig Wiener pe spectrograma zgomotoasa;
    \item reconstruieste forma de unda folosind aceleasi utilitare STFT/ISTFT ca restul proiectului.
\end{itemize}

Metoda este utila deoarece ofera o baza matematica simpla pentru enhancement audio. Spre deosebire de modelul neuronal, filtrul Wiener nu are nevoie de antrenare. El se bazeaza pe presupunerea ca zgomotul poate fi aproximat printr-un profil stationar, estimat din segmentele mai putin energice ale semnalului.

\chapter{Modelul matematic}

Se considera semnalul observat:

\[
y[n] = x[n] + d[n],
\]

unde \(x[n]\) este semnalul vocal curat, iar \(d[n]\) este zgomotul adaugat. In domeniul STFT, relatia devine:

\[
Y(f,t) = X(f,t) + D(f,t),
\]

unde \(f\) reprezinta binul de frecventa, iar \(t\) reprezinta cadrul temporal.

Filtrul Wiener construieste un castig \(G(f,t)\), aplicat pe magnitudinea spectrogramei zgomotoase:

\[
|\hat{X}(f,t)| = G(f,t) \cdot |Y(f,t)|.
\]

In implementarea din proiect, castigul este calculat din raportul semnal-zgomot posterior:

\[
\operatorname{SNR}_{post}(f,t) =
\frac{|Y(f,t)|^2}{P_D(f,t) + \varepsilon},
\]

unde \(P_D(f,t)\) este estimarea puterii zgomotului, iar \(\varepsilon\) previne impartirea la zero.

Castigul final este:

\[
G(f,t) =
\frac{\operatorname{SNR}_{post}(f,t)}
{\operatorname{SNR}_{post}(f,t) + 1}.
\]

Cand energia observata este mult mai mare decat zgomotul estimat, \(G\) se apropie de 1 si semnalul este pastrat. Cand energia este apropiata de zgomot, \(G\) scade si componenta respectiva este atenuata.

\chapter{Importuri si dependinte}

\section{Liniile 1--4}

\codechunk{1}{4}

\begin{description}[leftmargin=2.8cm,style=nextline]
    \item[Linia 1] Importa biblioteca \texttt{numpy}, folosita pentru operatii numerice pe spectrograme complexe, vectori de energie, medii si valori absolute.
    \item[Linia 2] Importa \texttt{torch}, deoarece pipeline-ul proiectului foloseste STFT si ISTFT bazate pe tensori PyTorch pentru consistenta cu modelul neuronal.
    \item[Linia 3] Importa modulul \texttt{stft\_utils}, care contine functii comune pentru calculul spectrogramei, separarea magnitudine-faza si reconstructia semnalului.
    \item[Linia 4] Importa \texttt{config}, de unde sunt preluati parametrii globali ai reprezentarii audio: \texttt{N\_FFT}, \texttt{HOP\_LEN} si \texttt{FRAME\_LEN}.
\end{description}

Aceste dependinte arata ca filtrul Wiener nu este implementat ca experiment izolat, ci ca metoda integrata in aceeasi conventie de procesare audio folosita de restul proiectului.

\chapter{Estimarea zgomotului}

\section{Liniile 6--13}

\codechunk{6}{13}

Functia \texttt{estimate\_noise\_psd} estimeaza profilul spectral al zgomotului. Ea primeste o spectrograma complexa \texttt{S\_noisy}, adica reprezentarea STFT a semnalului zgomotos, si intoarce o estimare a densitatii spectrale de putere a zgomotului.

\begin{description}[leftmargin=2.8cm,style=nextline]
    \item[Linia 6] Defineste functia. Parametrul \texttt{S\_noisy} este un tablou NumPy complex, iar \texttt{num\_noise\_frames} stabileste cate cadre cu energie mica vor fi folosite pentru estimarea zgomotului. Valoarea implicita este 6.
    \item[Linia 7] Documenteaza ideea functiei: zgomotul este estimat din cadrele cu cea mai mica energie, fara etichete VAD externe.
    \item[Linia 8] Apeleaza \texttt{stft\_utils.mag\_phase}. Se extrage magnitudinea spectrogramei, iar faza este ignorata deoarece estimarea puterii zgomotului nu are nevoie de faza.
    \item[Linia 9] Citeste numarul de cadre temporale din a doua dimensiune a matricei de magnitudine. Forma tipica este \((F,T)\), unde \(F\) este numarul de binuri de frecventa si \(T\) numarul de cadre.
    \item[Linia 10] Calculeaza energia medie a fiecarui cadru:
    \[
    E(t) = \frac{1}{F}\sum_f |Y(f,t)|^2.
    \]
    Rezultatul este un vector cu cate o valoare de energie pentru fiecare cadru temporal.
    \item[Linia 11] Sorteaza indicii cadrelor dupa energie si pastreaza primele cadre, adica cele mai slabe energetic. Expresia \texttt{min(num\_noise\_frames, num\_frames)} evita selectarea mai multor cadre decat exista in semnal.
    \item[Linia 12] Calculeaza media puterii pe cadrele selectate, pentru fiecare frecventa. Parametrul \texttt{keepdims=True} pastreaza rezultatul sub forma \((F,1)\), forma convenabila pentru difuzare pe toate cadrele.
    \item[Linia 13] Returneaza estimarea zgomotului.
\end{description}

Aceasta functie presupune ca in semnal exista cel putin cateva cadre in care vorbirea este absenta sau mai slaba. In practica, aceasta presupunere este frecventa in semnalele de vorbire, dar nu este garantata in toate inregistrarile. Daca vorbirea este continua si foarte intensa, estimarea zgomotului poate contine si energie vocala.

\chapter{Filtrarea Wiener in domeniul STFT}

\section{Liniile 15--27}

\codechunk{15}{27}

Functia \texttt{wiener\_filter\_stft} aplica filtrul Wiener pe o spectrograma complexa. Ea nu citeste direct un fisier audio si nu reconstruieste forma de unda; lucreaza strict in domeniul STFT, ceea ce o face reutilizabila in experimente si pipeline-uri.

\begin{description}[leftmargin=2.8cm,style=nextline]
    \item[Linia 15] Defineste functia. \texttt{S\_noisy} este spectrograma complexa de intrare, \texttt{noise\_psd} este estimarea puterii zgomotului, iar \texttt{eps} este o constanta mica pentru stabilitate numerica.
    \item[Liniile 16--18] Contin docstring-ul functiei. El precizeaza ca filtrarea foloseste un profil intern de zgomot stationar.
    \item[Linia 19] Separa spectrograma complexa in magnitudine si faza. Filtrul modifica magnitudinea, dar pastreaza faza initiala.
    \item[Liniile 20--21] Daca \texttt{noise\_psd} are o singura coloana, aceasta este repetata pe toate cadrele. Astfel, un profil de zgomot estimat global pentru fiecare frecventa devine o matrice compatibila cu spectrograma \((F,T)\).
    \item[Linia 22] Calculeaza puterea semnalului observat:
    \[
    |Y(f,t)|^2.
    \]
    \item[Linia 23] Calculeaza SNR-ul posterior prin raportarea puterii observate la puterea zgomotului estimat.
    \item[Linia 24] Calculeaza castigul Wiener. Formula produce valori intre 0 si 1 pentru valori nenegative ale SNR-ului.
    \item[Linia 25] Aplica filtrul asupra magnitudinii. Binurile dominate de zgomot sunt reduse, iar cele dominate de semnal sunt pastrate mai mult.
    \item[Linia 26] Reconstruieste spectrograma complexa imbunatatita din magnitudinea filtrata si faza originala.
    \item[Linia 27] Returneaza spectrograma complexa imbunatatita.
\end{description}

Un detaliu important este pastrarea fazei originale. Aceasta este o alegere comuna in metodele clasice de enhancement, deoarece estimarea fazei curate este dificila. In multe sisteme, imbunatatirea magnitudinii produce deja o crestere perceptibila a claritatii, chiar daca faza ramane cea zgomotoasa.

\chapter{Functia de procesare a formei de unda}

\section{Liniile 29--65}

\codechunk{29}{65}

Functia \texttt{enhance\_waveform} este interfata practica a modulului. Ea primeste un semnal audio in domeniul timp, il transforma in STFT, aplica estimarea de zgomot si filtrul Wiener, apoi revine la forma de unda.

\section{Pregatirea semnalului}

\begin{description}[leftmargin=2.8cm,style=nextline]
    \item[Linia 29] Defineste functia care primeste \texttt{noisy}, un semnal audio zgomotos sub forma de tablou NumPy.
    \item[Liniile 30--31] Daca semnalul are mai multe canale, se face media pe canale. Rezultatul este un semnal mono, compatibil cu restul pipeline-ului.
    \item[Linia 32] Converteste semnalul la \texttt{float32}. Aceasta este forma asteptata de majoritatea operatiilor audio si de PyTorch.
    \item[Linia 34] Comentariul explica motivul folosirii STFT-ului Torch: consistenta cu metodele bazate pe model.
    \item[Linia 35] Converteste tabloul NumPy in tensor PyTorch.
\end{description}

\section{Transformarea STFT}

\begin{description}[leftmargin=2.8cm,style=nextline]
    \item[Liniile 36--41] Apeleaza \texttt{stft\_utils.compute\_stft}. Parametrii sunt luati din \texttt{config}, ceea ce garanteaza ca metoda Wiener foloseste aceeasi rezolutie timp-frecventa ca restul proiectului.
    \item[Linia 38] \texttt{config.N\_FFT} controleaza numarul de puncte FFT.
    \item[Linia 39] \texttt{config.HOP\_LEN} controleaza pasul dintre cadre.
    \item[Linia 40] \texttt{config.FRAME\_LEN} controleaza lungimea ferestrei.
\end{description}

Rezultatul este impartit in doua componente: \texttt{noisy\_mag}, magnitudinea, si \texttt{noisy\_phase}, faza. Aceasta reprezentare este folosita frecvent in sistemele de reducere a zgomotului, deoarece multe metode modifica doar magnitudinea.

\section{Aplicarea filtrului}

\begin{description}[leftmargin=2.8cm,style=nextline]
    \item[Linia 43] Comentariul marcheaza conversia inapoi la o spectrograma complexa, necesara pentru functiile definite mai sus.
    \item[Linia 44] Reconstruieste spectrograma complexa:
    \[
    S_{noisy}(f,t) = |Y(f,t)|e^{j\angle Y(f,t)}.
    \]
    In cod, aceasta relatie este implementata prin inmultirea magnitudinii cu exponentiala complexa a fazei.
    \item[Linia 45] Estimeaza puterea zgomotului din spectrograma zgomotoasa.
    \item[Linia 46] Aplica filtrul Wiener si obtine spectrograma imbunatatita.
\end{description}

\section{Reconstructia audio}

\begin{description}[leftmargin=2.8cm,style=nextline]
    \item[Linia 48] Comentariul indica extragerea magnitudinii si fazei din spectrograma imbunatatita.
    \item[Linia 49] Calculeaza magnitudinea spectrogramei filtrate.
    \item[Linia 50] Calculeaza faza spectrogramei filtrate.
    \item[Linia 52] Comentariul explica revenirea la domeniul timp prin ISTFT Torch.
    \item[Liniile 53--54] Convertesc magnitudinea si faza inapoi in tensori PyTorch.
    \item[Liniile 55--62] Apeleaza \texttt{stft\_utils.inverse\_stft}, folosind aceiasi parametri STFT ca la analiza. Parametrul \texttt{length=len(noisy)} forteaza lungimea iesirii sa fie egala cu lungimea semnalului initial.
    \item[Linia 64] Converteste tensorul rezultat inapoi in NumPy.
    \item[Linia 65] Returneaza forma de unda imbunatatita.
\end{description}

\chapter{Fluxul complet al datelor}

Fluxul implementat de fisier poate fi rezumat astfel:

\begin{enumerate}
    \item intrarea este un semnal audio zgomotos in domeniul timp;
    \item semnalul este convertit la mono si \texttt{float32};
    \item se calculeaza STFT-ul folosind parametrii globali;
    \item magnitudinea si faza sunt combinate intr-o spectrograma complexa;
    \item zgomotul este estimat din cadrele cu energie minima;
    \item se calculeaza castigul Wiener;
    \item magnitudinea este atenuata in zonele dominate de zgomot;
    \item spectrograma filtrata este reconstruita;
    \item ISTFT-ul produce semnalul audio final.
\end{enumerate}

Acest flux face ca metoda sa poata fi comparata corect cu alte metode ale proiectului, deoarece foloseste aceleasi conventii de segmentare si reconstructie audio.

\chapter{Avantaje, limitari si observatii}

\section{Avantaje}

\begin{itemize}
    \item Nu necesita antrenare si poate fi rulat direct pe orice semnal.
    \item Este rapid, fiind bazat pe operatii vectorizate NumPy si PyTorch.
    \item Este usor de interpretat matematic.
    \item Este compatibil cu pipeline-ul STFT al proiectului.
    \item Poate servi ca baseline pentru metode neuronale mai complexe.
\end{itemize}

\section{Limitari}

\begin{itemize}
    \item Presupune ca zgomotul este relativ stationar.
    \item Estimarea zgomotului poate fi gresita daca toate cadrele contin vorbire puternica.
    \item Pastreaza faza zgomotoasa, ceea ce limiteaza calitatea reconstructiei.
    \item Poate produce atenuare excesiva in frecventele unde vocea are energie apropiata de zgomot.
    \item Nu invata structuri vocale complexe, spre deosebire de modelul neuronal.
\end{itemize}

\chapter{Concluzie}

Fisierul \texttt{wiener\_filter.py} implementeaza o metoda DSP clasica si importanta pentru reducerea zgomotului. El are rolul de metoda de referinta in proiect: este simplu, explicabil si integrat in acelasi pipeline audio ca metodele neuronale.

Prin functiile sale, fisierul acopera intregul traseu de procesare: estimarea zgomotului, calculul castigului Wiener, filtrarea spectrogramei si reconstructia semnalului. Din acest motiv, este unul dintre fisierele importante ale proiectului atunci cand se explica partea de procesare audio clasica si comparatia dintre metode deterministe si metode bazate pe invatare automata.

\end{document}
